% From Nothing to Coulomb's Law — lesson notes (12 July 2026)
% Build:  pdflatex -jobname=2026-07-12-coulomb-law-notes coulomb-law-notes.tex
\documentclass[11pt,a4paper]{article}
\usepackage[margin=2.2cm,bottom=2.4cm]{geometry}
\usepackage[T1]{fontenc}
\usepackage[utf8]{inputenc}
\usepackage{lmodern}
\usepackage{amsmath}
\usepackage{xcolor}
\usepackage{booktabs}
\usepackage{enumitem}
\usepackage{parskip}
\usepackage{fancyhdr}
\usepackage{tcolorbox}
\usepackage{url}

\definecolor{accent}{HTML}{2F6FED}
\definecolor{muted}{HTML}{5B6473}
\definecolor{ink}{HTML}{1F2430}
\definecolor{boxbg}{HTML}{F2F6FF}
\definecolor{rule}{HTML}{E2E5EA}

\pagestyle{fancy}
\fancyhf{}
\renewcommand{\headrulewidth}{0pt}
\renewcommand{\footrulewidth}{0.4pt}
\fancyfoot[L]{\small\color{muted} From Nothing to Coulomb's Law --- lesson notes}
\fancyfoot[R]{\small\color{muted} Page \thepage}

\newcommand{\nsec}[1]{\vspace{8pt}\par{\color{accent}\large\bfseries #1}\par\vspace{2pt}}

\begin{document}\color{ink}

{\huge\bfseries From Nothing to Coulomb's Law}\\[4pt]
{\color{muted} Electricity --- lesson notes \;\textperiodcentered\; discussion of 12 July 2026 \;\textperiodcentered\; Dad's notes, corrected and extended}\\[6pt]
{\color{rule}\rule{\linewidth}{0.9pt}}

\nsec{The chain of ideas}
\begin{enumerate}[label=\textcolor{accent}{\textbf{\arabic*.}}, leftmargin=1.9em, itemsep=2.5pt, topsep=2pt]
  \item \textbf{Everything is made of particles.}\\
        {\small\color{muted} Matter is made of atoms, and atoms are made of protons, neutrons and electrons.}
  \item \textbf{Electric charge is a property of a particle.}\\
        {\small\color{muted} Just as a particle has mass, it can carry charge. Charge is measured in coulombs (C).}
  \item \textbf{Particles can be positive, negative, or neutral.}\\
        {\small\color{muted} The proton is positive, the electron is negative, the neutron is neutral.}
  \item \textbf{Charge is quantised: it comes in steps of the elementary charge $e$.}\\
        {\small\color{muted} Every observed charge is a whole-number multiple of $e$. The proton carries exactly $+e$
        and the electron exactly $-e$. (Fine print: quarks carry $\pm e/3$ and $\pm 2e/3$, but they are never seen alone.)}
  \item \textbf{Charged particles interact: they attract or repel each other.}\\
        {\small\color{muted} Neutral particles feel no electric force.}
  \item \textbf{Like charges repel; unlike charges attract.}\\
        {\small\color{muted} $+$ and $+$ push apart, $-$ and $-$ push apart, $+$ and $-$ pull together.}
  \item \textbf{The net charge of a body counts its imbalance of particles: $Q = (N_{+} - N_{-})\,e$.}\\
        {\small\color{muted} $N_{+}$ = number of positive charges (protons), $N_{-}$ = number of negative ones
        (electrons). Most bodies have $N_{+} \approx N_{-}$, so they are neutral.}
  \item \textbf{Charge is conserved.}\\
        {\small\color{muted} Charge is never created or destroyed --- it only moves from one body to another
        (e.g.\ rubbing a balloon moves electrons onto it).}
  \item \textbf{Coulomb's law tells us how strong the force is.}\\
        {\small\color{muted} See the box below --- the climax of the chain.}
\end{enumerate}

\begin{tcolorbox}[colback=boxbg, colframe=accent, boxrule=0.9pt, arc=2.5mm,
                  left=10pt, right=10pt, top=7pt, bottom=7pt,
                  title={\large\bfseries Coulomb's law}, coltitle=white,
                  before skip=14pt]
\[
  F \;=\; k\,\frac{|q_1|\,|q_2|}{r^{2}}
  \qquad\text{with}\qquad
  k \;=\; \frac{1}{4\pi\varepsilon_0} \;\approx\; 8.99\times 10^{9}\ \mathrm{N\,m^{2}/C^{2}}
\]
\begin{itemize}[leftmargin=1.4em, itemsep=1.5pt, topsep=3pt]
  \item $q_1$, $q_2$ --- the two charges (in coulombs); $r$ --- the distance between them (in metres).
        Valid for \textbf{point charges} (small compared with $r$).
  \item The force acts \textbf{along the line joining the charges}: like charges are pushed apart,
        unlike charges pulled together. Both charges feel the \textbf{same} magnitude of force, in
        opposite directions (Newton's third law).
  \item \textbf{Inverse square:} double the distance $\Rightarrow$ one quarter of the force.
\end{itemize}
\smallskip
\textbf{Worked example.} Two charges of $1\,\mu$C each, $10\ \mathrm{cm}$ apart:
\[
  F \;=\; 8.99\times 10^{9} \cdot \frac{\left(10^{-6}\right)^{2}}{(0.1)^{2}}
  \;\approx\; 0.9\ \mathrm{N}
  \quad\text{--- about the weight of a small apple.}
\]
{\small\color{muted} Scale check: two charges of $1$ C at $1$ m would push with $9\times10^{9}$ N
($\approx$ the weight of 900\,000 tonnes) --- that is why real charges are microcoulombs.}
\end{tcolorbox}

% keep the heading together with its table
\par\noindent\begin{minipage}{\linewidth}
\nsec{Constants worth remembering}
\renewcommand{\arraystretch}{1.25}
\begin{tabular}{@{}llll@{}}
\toprule
\textbf{Quantity} & \textbf{Symbol} & \textbf{Value} & \textbf{Note}\\
\midrule
Elementary charge    & $e$    & $1.602\times10^{-19}$ C  & {\small\color{muted} exact: $1.602\,176\,634\times10^{-19}$ C}\\
Electron charge      & $-e$   & $-1.602\times10^{-19}$ C & {\small\color{muted} proton: $+e$}\\
Coulomb constant     & $k = 1/(4\pi\varepsilon_0)$ & $8.99\times10^{9}\ \mathrm{N\,m^{2}/C^{2}}$ & {\small\color{muted} remember as $\approx 9\times10^{9}$}\\
Vacuum permittivity  & $\varepsilon_0$ & $8.854\times10^{-12}\ \mathrm{C^{2}/(N\,m^{2})}$ & {\small\color{muted} the same thing as F/m}\\
Charges in 1 coulomb & $1\,\mathrm{C}/e$ & $6.24\times10^{18}$ & {\small\color{muted} 1 C is a \emph{huge} charge}\\
Electron mass        & $m_e$  & $9.11\times10^{-31}$ kg  & {\small\color{muted} proton is ${\approx}1836\times$ heavier}\\
\bottomrule
\end{tabular}
\end{minipage}

\nsec{Where to read (from our archive)}
\begin{itemize}[leftmargin=1.2em, itemsep=3pt, topsep=2pt]
  \item \textbf{Siyavula, Physical Sciences Gr 10} ---
        {\small\path{docs/textbook/siyavula-physical-sciences/Gr10_PhysicalSciences_Learner_Eng.pdf}}\\
        {\small ch~16 \emph{Electrostatics} --- \S16.2 Two kinds of charge (PDF p.\,276),
        \S16.3 Conservation of charge (PDF p.\,282), \S16.4 Quantisation of charge (PDF p.\,284).
        Closest match to these notes.}
  \item \textbf{CK-12, Physical Science Workbook} ---
        {\small\path{docs/textbook/ck12-physical-science/CK-12-Physical-Science-MS-Workbook-with-Answers.pdf}}\\
        {\small lesson 23.1 \emph{Electric Charge} (printed p.\,276 $\approx$ PDF p.\,283) ---
        practice worksheets; answers in ch~26.23 (printed p.\,417).}
  \item \textbf{NCERT, Class 9 Science} ---
        {\small\path{docs/textbook/ncert-class9-physics/other-chapters/ch08-journey-inside-the-atom.pdf}}\\
        {\small \S8.2--8.6 (PDF pp.\,2--12) --- protons \& electrons, atomic number: why atoms are
        neutral and where the charges sit.}
  \item \textbf{Cambridge Lower Secondary, Stage 7} ---
        {\small\path{docs/textbook/cambridge/Cambridge_Lower_Secondary_Science_2ed_7_Learners_book.pdf}}\\
        {\small Unit 9, \S9.1 \emph{Flow of electricity} ($\approx$ p.\,287) --- moving charge $=$ current:
        the next step after these notes.}
\end{itemize}

{\small\color{muted} Full quantitative treatment: \textbf{Siyavula Grade 11} (free, openly licensed) --- easy to add to the archive.}

\end{document}
